This work presented \sensecv{}, an egocentric visual--inertial dataset for indoor assistive navigation. It contains 592 events collected across three locations, with videos, inertial signals, and metadata recorded by an open-source Android application. The protocol includes free-path, lateral-deviation, and stopping events, while accelerometer and gyroscope data support temporal delimitation and label verification.

With Location C reserved for testing, MobileNetV4 achieved the best performance and had the fewest parameters, reaching 95.58\% accuracy for path status and 89.97\% for deviation direction. The higher variability in deviation classification indicates greater difficulty generalizing this task to unseen environments.

Reducing sampling from 30 to 10~FPS preserved path-status performance and may improve deviation classification when temporal separation is maintained. Conversely, random frame-level splits produced artificially high results, confirming that correlated frames can inflate performance estimates.

Future work will increase the diversity of locations, participants, and devices, evaluate temporal models, and explore the preserved location and magnetometer signals.
